% SYN-ACK WIRE -- user manual for the Atari 8-bit computers
% Build: make   (or make in docs/: the rulebook and every manual)
%
% The style, the theme macros and the furniture the shared rules text needs are in ../../style/synack.sty;
% the rules text is ../../rules/*.tex; the Atari's names for things are in platform.tex.
% Screens are in screens/: real captures, and dumps drawn by client/tools/render_screen.py.
\documentclass[12pt,letterpaper]{article}

\usepackage{synack}
% The Atari's names for things, for the shared style and text (../../style/synack.sty, ../../rules/).
\renewcommand{\DocKind}{USER MANUAL // ATARI 8-BIT}
\renewcommand{\PlatformShort}{Atari}                       % as in "on the Atari"
\newcommand{\PlatformLong}{Atari 8-bit computers}          % the title page
\newcommand{\PlatformMachines}{ATARI 400 / 800 / XL / XE}  % the title page, the screen chapter
\newcommand{\SignalLost}{the title bar reads \code{NO SIGNAL} and turns purple}  % the link is quiet

% Where the 40 x 24 screen sits in a capture, in image pixels (see \capgeom):
\makeatletter
\newcommand{\rendered}{\capgeom{10}{8}{8}{8}{340}{208}}     % client/tools/render_screen.py
\newcommand{\emulated}{\capgeom{11}{11}{13.7}{16}{570}{408}} % synack_title_screen.png: an emulator window
\makeatother

\graphicspath{{screens/}}

\makeatletter
% The rules' captures are staged games (server/stage.py, scenarios.toml) on the board of the drawing
% under "1. Spike a node", so the drawing, the text and the screen tell one story.
\expandafter\def\csname shot@spike\endcsname{\screen{spike.png}{CAPTURE \quad \MakeUppercase{\V{spike}} A \MakeUppercase{\V{node}}}%
  {The drawing above, as ZER0COOL sees it. The move hints \code{1}--\code{4} sit on its four targets: they are the
  \Vp{exploit} in your \V{kit} that run there, as the info line lists them. Press the one to run.}}
\expandafter\def\csname shot@broker\endcsname{\screen{infiltration.png}{CAPTURE \quad \MakeUppercase{\V{infiltration}}}%
  {Target 2: ZER0COOL ran \node{C6}, next to the lone \V{spike} on \node{B6}. The info line lists, inverse, the
  syndicates with no \V{hive} on the \V{grid}: press one of those letters to \V{broker} the new \V{hive}.}}
\expandafter\def\csname shot@hivejack\endcsname{\screen{tied-root.png}{CAPTURE \quad \MakeUppercase{\V{hivejack}}: A TIED \MakeUppercase{\V{root}}}%
  {The example below, as ZER0COOL sees it. The spike on \node{G3} joined \Syn{5}{name} (4 \Vp{node}) and \Syn{3}{name}
  (3); the log shows the tied \V{root} paid. The info line says how many of the \V{absorbed} \Vp{coin} you hold: type
  how many to \V{cashout} and press \key{Return}, then how many to \V{swap}, and \key{Return} again.}}
\expandafter\def\csname shot@buy\endcsname{\screen{spread-buy.png}{CAPTURE \quad BUY \MakeUppercase{\V{coin}}}%
  {Target 3: \node{E1} spread \Syn{5}{name} to 5 \Vp{node}, and its price is already up. While you buy, the
  \Vp{hive} panel shows each price and how many the \V{bank} has left. Type up to \Rule{maxbuy} letters, in any
  mix (\code{\Syn{5}{letter}\Syn{5}{letter}\Syn{3}{letter}} buys two \Syn{5}{name} and one \Syn{3}{name}),
  \key{Delete} to take one back, \key{Return} to buy.}}
\expandafter\def\csname shot@nullroute\endcsname{\screen{nullroute.png}{CAPTURE \quad \MakeUppercase{\V{lockdown}}: THE QUESTION}%
  {ZER0COOL's \V{spike} on \node{H2} spread \Syn{2}{name} to \Rule{safesize} \Vp{node}: \V{safe}, and with
  \Syn{5}{name} and \Syn{0}{name} already \V{safe}, every \V{hive} is. With too little \Game{currency} for any
  \V{coin}, the prompt is the question alone: \key{Y} calls the \V{nullroute}, \key{N} or \key{Return} plays on.}}
\expandafter\def\csname shot@lobby\endcsname{\screen{lobby-join.png}{CAPTURE \quad JOINING A LOBBY}%
  {ZER0COOL hosts and MAMBA is waiting; the Atari has just joined. The host clock started when MAMBA arrived, and every
  newcomer is told how long the host has left.}}
\expandafter\def\csname shot@final\endcsname{\screen{game-over.png}{CAPTURE \quad THE END OF A GAME}%
  {ZER0COOL called it. \VC{root} and \V{sudo} were paid for every \V{hive}, all \V{coin} was sold, and the log ends
  with each \V{player}'s final \Game{currency} and the top \V{player}. The table is a lobby again: the host's prompt
  asks for the seats of a rematch.}}
\makeatother


\begin{document}

% ================================================================== title
\begin{titlepage}
\machine
\vspace*{1.6in}
{\color{amber}\rule{\linewidth}{2pt}}\par\vspace{14pt}
{\fontsize{44}{48}\selectfont \game\par}\vspace{10pt}
{\fontsize{12}{16}\selectfont\color{steel} SEVEN SYNDICATES. ONE REGION. EVERY \MakeUppercase{\V{node}} FOR SALE.\par}
\vspace{14pt}{\color{amber}\rule{\linewidth}{2pt}}\par
\vspace{0.35in}
{\fontsize{24}{30}\selectfont USER MANUAL\par}\vspace{4pt}
{\fontsize{14}{18}\selectfont\color{amber} FOR THE \MakeUppercase{\PlatformLong}\par}
\vspace{0.4in}
{\fontsize{12}{16}\selectfont \PlatformMachines \quad // \quad FUJINET\par
\Rule{minplayers}-\Rule{maxplayers} \MakeUppercase{\Vp{player}} \quad // \quad BOTS FILL EMPTY SEATS\par}
\vfill
{\fontsize{8.5}{11}\selectfont\color{steel}
AN UNOFFICIAL CYBERPUNK RE-SKIN OF ACQUIRE BY SID SACKSON.\par
ACQUIRE IS A TRADEMARK OF HASBRO, INC. NOT AFFILIATED WITH OR ENDORSED BY HASBRO.\par}
\end{titlepage}

% ================================================================== 01
\renewcommand{\secstatus}{THE SETTING \quad // \quad THE OBJECT \quad // \quad THE END}
\chapter{The Network Grid}
\begin{multicols}{2}
\rules{opening}
\end{multicols}

% ================================================================== 02
\renewcommand{\secstatus}{REGION \quad // \quad STATUS \quad // \quad HANDLE}
\chapter{\VC{jackin}}
Boot the \game\ disk, or pick a \game\ server in the FujiNet lobby. The lobby passes the server you pick through FujiNet's SD card: without one, you join CHIBA-GRID-7. The title screen dials the server as
soon as it appears. Type a handle, up to 8 characters, and press \key{Return}.

\begin{capture}[\emulated]{synack_title_screen.png}{CAPTURE 01 \quad TITLE SCREEN \quad LOBBY OPEN}
  \callout{1}{L}{4}{16}
  \callout{2}{R}{33}{17}
  \callout{3}{L}{4}{18}
  \callout{4}{L}{4}{20}
  \callout{5}{R}{34}{21}
  \callout{6}{R}{30}{23}
\end{capture}

\begin{legend}
\li{1}{\textbf{Region.} The server's \V{grid}: one game at a time, up to six seats.}
\li{2}{\textbf{Status.} Whether you can join, and how: \code{LOBBY OPEN}, \code{IN PROGRESS}, \code{TABLE FULL}.
It stays up to date while you type.}
\li{3}{\textbf{Handle.} Your name at the table. From the FujiNet lobby, it starts with your lobby name.}
\li{4}{\key{Option} turns the move hints on or off.}
\li{5}{\key{?} opens the instructions, here or in the middle of a game. On an XL or XE, \key{Help} does too.}
\li{6}{\textbf{The prompt line.} What the screen expects from you next.}
\end{legend}

% ================================================================== 03
\renewcommand{\secstatus}{\PlatformMachines \quad // \quad GR.0 TEXT \quad // \quad 40 x 24}
\chapter{The Screen}
Everything you need is on one screen. This is early in a game, and it's your turn: pick an \V{exploit}
to run.

\begin{capture}{spike-early.png}{CAPTURE 02 \quad \MakeUppercase{\V{spike}} PROMPT \quad 3 SEATS \quad EARLY GAME}
  \callout{1}{L}{2}{0}
  \callout{2}{L}{2}{6}
  \callout{3}{R}{35}{1}
  \callout{4}{R}{23}{10}
  \callout{5}{L}{0}{13}
  \callout{6}{L}{1}{20}
  \callout{7}{L}{1}{22}
  \callout{8}{R}{38}{23}
\end{capture}

\begin{legend}
\li{1}{\textbf{Title bar.} The region and the link. It reads \code{NO SIGNAL} when the server goes quiet.}
\li{2}{\textbf{The \V{grid}.} Inverse letters are \Vp{hive}; \node{\Glyph{spike}} is a lone \V{spike}. A digit marks
a \V{node} an \V{exploit} in your \V{kit} can run on. The last move blinks for a few seconds.}
\li{3}{\textbf{\VCp{hive}.} Each syndicate's letter, \V{hive} size and name. While you buy, its price and how many
the \V{bank} has left.}
\li{4}{\textbf{Last move.} The \V{node} \V{spike}d most recently.}
\li{5}{\textbf{\VCp{player}.} \code{>} marks you, \code{*} a bot, \code{?} someone away. Account, then
\V{coin} held. An inverse name has the turn.}
\li{6}{\textbf{The log.} The last three things that happened at the table.}
\li{7}{\textbf{Your \V{kit}.} Your \Rule{kitsize} \Vp{exploit} and the \Vp{node} they run on.}
\li{8}{\textbf{The prompt.} What you're asked, with seconds left on the shot clock. Purple means the
decision is yours. Here: press \key{1}--\key{6}.}
\end{legend}

% ================================================================== 04
\renewcommand{\secstatus}{THE KEYBOARD \quad // \quad CONSOLE KEYS \quad // \quad NO JOYSTICK}
\chapter{Keys}
\game\ is played from the keyboard. Each question on the prompt line names the keys it takes; this chapter
has them all in one place. The game checks an answer before sending it and says on the prompt line what's
wrong, such as \code{NOT ENOUGH NEW YEN}.

\subsection{Any time}
\begin{keytable}
\kr{\key{?}\ \ \key{Help}}{Open the instructions, on the title screen or in the middle of a game. \key{Help} is on the XL and XE.}
\kr{\key{Option}}{Move hints on or off: the digits on the \V{grid} that show where each \V{exploit} in your \V{kit} can run.}
\kr{\key{Select}}{Sound on or off. A chord tells you it's on.}
\end{keytable}

\subsection{The title screen}
\begin{keytable}
\kr{\key{A}--\key{Z}\ \ \key{0}--\key{9}}{Type your handle, up to 8 characters.}
\kr{\key{Delete}}{Rub out the last character.}
\kr{\key{Return}}{\VC{jackin}.}
\end{keytable}

\subsection{The instructions}
\begin{keytable}
\kr{\key{>}\ \ \key{\fontsize{12}{9}\selectfont$\rightarrow$}}{Next page.}
\kr{\key{<}\ \ \key{\fontsize{12}{9}\selectfont$\leftarrow$}}{Previous page.}
\kr{\key{Esc}}{Back to where you were. \key{Return}, \key{Help} and \key{?} do the same.}
\end{keytable}

\subsection{Hosting}
\begin{keytable}
\kr{\key{2}--\key{6}}{Start the game with that many seats. Bots take the seats nobody has joined.}
\end{keytable}

\subsection{Your turn}
\begin{keytable}
\kr{\key{1}--\key{6}}{\VC{spike}: run that \V{exploit} from your \V{kit}. An inverse one can't run now; press it and the prompt line says why.}
\kr{\key{A}--\key{Z}}{\VC{infiltration}: the letter of the syndicate to \V{broker} the new \V{hive} to.}
\kr{\key{A}--\key{Z}}{A tied \V{hivejack}: the letter of the \V{hive} that survives, or is \V{absorbed} next.}
\kr{\key{A}--\key{Z}}{Buy \V{coin}: type the syndicate letters, up to \Rule{maxbuy}, in any mix of the \Vp{hive} on the \V{grid}.
One letter per \V{coin}: \key{\Syn{5}{letter}} buys one, \key{\Syn{5}{letter}}\key{\Syn{5}{letter}} two, \key{\Syn{5}{letter}}\key{\Syn{5}{letter}}\key{\Syn{5}{letter}} three of the same, \key{\Syn{5}{letter}}\key{\Syn{0}{letter}} one each of two, \key{\Syn{5}{letter}}\key{\Syn{0}{letter}}\key{\Syn{6}{letter}} one each of three.}
\kr{\key{Delete}}{Take back the last \V{coin} in a buy.}
\kr{\key{Return}}{Confirm the buy. \key{Return} alone buys nothing.}
\kr{\key{Y}\ \ \key{N}}{When the game may end: \V{nullroute} the region now, or play on. \key{Return} plays on too.}
\end{keytable}

\subsection{When your \V{hive} is \V{absorbed}}
\begin{keytable}
\kr{\key{0}--\key{9}\ \ \key{Return}}{First, how many \Vp{coin} to \V{cashout}; then how many to \V{swap}, an even number.
The rest you \V{hodl}. \key{Delete} rubs out a digit.}
\end{keytable}

\subsection{Leaving}
\begin{keytable}
\kr{\key{Esc}\ \ \key{Y}}{\VC{logoff} for good. A bot takes over your seat at once, and you're back on the title screen.
Any key but \key{Y} stays.}
\end{keytable}

% ================================================================== 05
\renewcommand{\secstatus}{THE \MakeUppercase{\V{grid}} \quad // \quad THE SYNDICATES \quad // \quad SETUP}
\chapter{What's in the Game}
\begin{multicols}{2}
\rules{contents}
\end{multicols}

\subsection{The Seven Syndicates}
\syntable

\subsection{Setup}
\begin{multicols}{2}
\rules{setup}
\end{multicols}

% ================================================================== 06
\renewcommand{\secstatus}{\MakeUppercase{\V{spike}} \quad // \quad BUY \quad // \quad DRAW}
\chapter{Playing a Turn}
\begin{multicols}{2}
\rules{turn}
\end{multicols}

% ================================================================== 07
\renewcommand{\secstatus}{\MakeUppercase{\V{totalcompromise}} \quad // \quad \MakeUppercase{\V{lockdown}} \quad // \quad \MakeUppercase{\V{cashout}}}
\chapter{Ending the Game}
\begin{multicols}{2}
\rules{ending}
\end{multicols}

% ================================================================== 08
\renewcommand{\secstatus}{THE LOBBY \quad // \quad THE HOST \quad // \quad SEATS AND BOTS}
\chapter{Hosting, Joining and Leaving}
\begin{multicols}{2}
\rules{hosting}
\end{multicols}

% ================================================================== 09
\renewcommand{\secstatus}{\MakeUppercase{\V{finderscut}} \quad // \quad \MakeUppercase{\V{root}} AND \MakeUppercase{\V{sudo}} \quad // \quad CASH}
\chapter{Where the Money Comes From}
\begin{multicols}{2}
\rules{money}
\end{multicols}

% ================================================================== 10
\renewcommand{\secstatus}{EARLY \quad // \quad MIDDLE \quad // \quad LATE}
\vspace{24pt}\section{Tips}                  % follows on: the money chapter ends early on its page
\begin{multicols}{2}
\rules{tips}
\end{multicols}

% ================================================================== 11
\renewcommand{\secstatus}{THE RULES \quad // \quad THE ATARI}
\chapter{Frequently Asked Questions}
\begin{multicols}{2}
\rules{faq}

\subsection{On the Atari}
\begin{description}[leftmargin=0pt,labelsep=0.5em]
  \item[Can I play alone?] Yes. Host a game and press \key{\Rule{minplayers}}: one bot plays you. Any seat nobody
    takes goes to a bot.
  \item[What are the digits on the \V{grid}?] Move hints: each marks a \V{node} that an \V{exploit} in your \V{kit}
    can run on, and the digit is its key. \key{Option} turns them on or off.
  \item[What do \code{*} and \code{?} after a name mean?] \code{*} is a bot. \code{?} is a \V{player} who is away,
    or on autopilot; a bot plays for them meanwhile.
  \item[Can I turn the sound off?] Yes: \key{Select}, any time. A chord says it's back on.
  \item[Why was my handle refused?] Someone at the table has it, a bot included. If it's your own seat after a
    reboot, the server may not have noticed the drop yet: try again in a minute.
  \item[What if everybody leaves?] The game waits for 10 minutes with nobody connected, then it is dropped.
\end{description}
\end{multicols}

% ================================================================== 12
\renewcommand{\secstatus}{EVERY TERM, FROM THE THEME}
\chapter{Glossary}
\rules{glossary}

% ================================================================== REF
\renewcommand{\secstatus}{FIND THE \MakeUppercase{\V{hive}}'S SIZE IN ITS TIER, THEN READ ACROSS}
\clearpage\renewcommand{\secnum}{REF}% after the page break: the running head uses it
\section{\V{coin} Prices and Payouts}
\vspace{6pt}
\begin{center}
\def\textcolor#1#2{\inv{#2}}% the tier's syndicate letters inverse, as on the screen
\newcommand{\hot}[1]{\leavevmode{\color{alarm}#1}}
\machine\fontsize{12}{14}\selectfont
\renewcommand{\arraystretch}{1.3}
\newcolumntype{C}{>{\centering\arraybackslash}p{0.85in}}
\newcolumntype{R}{>{\raggedleft\arraybackslash}p{0.95in}}
\begin{tabular}{CCC@{\hspace{14pt}}R@{\hspace{14pt}}RR}
LOW & MID & HIGH & \leavevmode\color{steel}PRICE & \leavevmode\color{signal}\MakeUppercase{\V{root}} & \leavevmode\color{signal}\MakeUppercase{\V{sudo}} \\
\TierCoins{low} & \TierCoins{mid} & \TierCoins{high} & & & \\[2pt]
\hline
2      & -      & -      & \Y{200}   & \Y{2,000}  & \Y{1,000} \\
3      & 2      & -      & \Y{300}   & \Y{3,000}  & \Y{1,500} \\
4      & 3      & 2      & \Y{400}   & \Y{4,000}  & \Y{2,000} \\
5      & 4      & 3      & \Y{500}   & \Y{5,000}  & \Y{2,500} \\
6-10   & 5      & 4      & \Y{600}   & \Y{6,000}  & \Y{3,000} \\
\hot{11-20}  & 6-10   & 5      & \Y{700}   & \Y{7,000}  & \Y{3,500} \\
21-30  & \hot{11-20}  & 6-10   & \Y{800}   & \Y{8,000}  & \Y{4,000} \\
31-40  & 21-30  & \hot{11-20}  & \Y{900}   & \Y{9,000}  & \Y{4,500} \\
41+    & 31-40  & 21-30  & \Y{1,000} & \Y{10,000} & \Y{5,000} \\
-      & 41+    & 31-40  & \Y{1,100} & \Y{11,000} & \Y{5,500} \\
-      & -      & 41+    & \Y{1,200} & \Y{12,000} & \Y{6,000} \\
\hline
\end{tabular}
\end{center}
\vspace{10pt}
\begin{center}\machine\fontsize{10}{16}\selectfont
\begin{tabular}{@{}ll@{}}
\textcolor{alarm}{RED} & FIRST \MakeUppercase{\V{safe}} SIZE: THE \MakeUppercase{\V{hive}} CAN'T BE \MakeUppercase{\V{absorbed}}\\
\textcolor{signal}{\MakeUppercase{\V{root}}} & 10 x \MakeUppercase{\V{coin}} PRICE\\
\textcolor{signal}{\MakeUppercase{\V{sudo}}} & 5 x \MakeUppercase{\V{coin}} PRICE\\
\Rule{endsize}+ & \MakeUppercase{\V{totalcompromise}}: THE \MakeUppercase{\V{nullroute}} MAY BE CALLED
\end{tabular}\end{center}

\vspace{6pt}
\begin{center}
{\machine\fontsize{9}{11}\selectfont\color{steel} THE LETTERS: EACH SYNDICATE ON THE \MakeUppercase{\V{grid}}, AND ITS TIER}\par\vspace{6pt}
\renewcommand{\arraystretch}{1.3}
\newcommand{\synrow}[1]{\inv{\Syn{#1}{letter}} & \Syn{#1}{name} & {\machine\color{steel}\MakeUppercase{\Syn{#1}{tier}}}}
\begin{tabular}{@{}cl@{\hspace{8pt}}l@{\hspace{36pt}}cl@{\hspace{8pt}}l@{}}
\synrow{0} & \synrow{4} \\
\synrow{1} & \synrow{5} \\
\synrow{2} & \synrow{6} \\
\synrow{3} & & & \\
\end{tabular}
\end{center}

\end{document}
